\section{About This Book}\label{sec:scope}

This is a companion for keeping Advent with the Roman Missal, the
\textit{Lectionary for Mass} and the \textit{Liturgy of the Hours} as they are
used in the dioceses of the United States. It follows the season from its first
evening to the night of Christmas. It is not a liturgical book or a guide to
the law: for the services of your own church and for what the Church requires
of you, follow your parish and its calendar. Statements about the calendar and
the law are as they stand at Advent 2026.

\subsection{The books it uses}

\textit{The Mass.} The Mass is that of the Roman Missal, third typical edition
(in Latin, 2002, emended 2008), in the English edition for the dioceses of the
United States (2011). The prayers of that Missal are described and not printed,
because their English is not free to reprint and their Latin, where it is the
work of the reform, is not printed either. Where an older Roman prayer stands
behind a present one, the older prayer is quoted as itself, from the Veronese
(or Leonine) Sacramentary, the Gelasian Sacramentary and the Gregorian
Sacramentary sent by Pope Hadrian~I to Charlemagne, and the book calls it the
older prayer, not the prayer of the Missal. The Entrance Antiphon of the Fourth
Sunday, whose words are those of the Introit of the Missal before 1970, is given
in the Latin of the Roman Missal printed at Tours in 1922 and the English of
F.~X. Lasance's \textit{New Roman Missal} (1945). The prayer of the Angelus,
whose words are the body of the Collect of the Fourth Sunday, is given in Latin
and English as the Raccolta of 1910 prints it.

\textit{The readings.} The readings of Mass are those of the \textit{Lectionary
for Mass for Use in the Dioceses of the United States}, second typical edition.
They are cited and not printed, by the Lectionary's own chapters and verses,
which are those of the New American Bible, with the Psalms numbered as in the
Hebrew. Where the Lectionary of the United States differs from the Latin
\textit{Ordo lectionum Missae} of 1981, the book follows the Lectionary of the
United States.

\textit{The Office.} The Office is that of the \textit{Liturgy of the Hours}.
Its second readings in the Office of Readings are named by author and work, and
quoted, where they are quoted, from the older English translations listed in
the References; the present English of the Hours is not printed. Where the
present Office keeps a text of the older Roman Breviary, the Latin is that of
the Breviary printed at Ratisbon in 1888 and the English that of John, Marquess
of Bute, in the edition of 1908.

\textit{Scripture.} Scripture in English is the Douay-Rheims Bible in Bishop
Challoner's revision, and Scripture in Latin is the Clementine Vulgate.
Quotations are cited by the Douay's chapter and verse under the usual name of
the book; where the Douay's name or numbering differs from the Lectionary's
(the Douay's 2~Kings is 2~Samuel, and its Ecclesiasticus is Sirach), the
difference is noted.

\textit{The Fathers and other writers.} Most of the Fathers are quoted from the
Ante-Nicene Fathers and the Nicene and Post-Nicene Fathers. St Bernard is quoted
in the Dublin translation of his sermons for the seasons (1921) and the London
translation of his sermons on Advent and Christmas (1909); St Anselm in Sidney
Norton Deane's translation (1903); St Bonaventure's life of St Francis in the
translation of 1904; St Francis Xavier in Henry James Coleridge's \textit{Life
and Letters} (1872); Blessed Pius IX's \latin{Ineffabilis Deus} in the bilingual
edition of 1855; St Thérèse of Lisieux in T.~N. Taylor's English of her
autobiography (1912); the councils in the English translation of Hefele's
history; and St Gregory of Tours in Dalton's translation. The other editions
are named in the References. Pope Francis's letter on the crib,
\latin{Admirabile signum}, and the other documents of the Holy See issued
after 1975 are given in the book's own words, with their numbers.

\textit{The Catena Aurea.} Some short sayings of the Fathers on the Gospels
come from the \textit{Catena Aurea}, the ``golden chain'' in which St Thomas
Aquinas set the Fathers' comments verse by verse beside the Gospels, quoted in
the Oxford translation. The Catena sometimes shortens or joins its sources, and
each saying is set under the Father to whom the Catena gives it.

\textit{The book's own translations.} Short renderings of Latin in the text,
given after the Latin they translate, are the author's own. So are the English
of the older Roman prayers behind the Collects and Prefaces of the Missal; of
the Roman calendar of 354; of the Roman Martyrology of 1902 on St Nicholas, St
Ambrose, St Damasus and St Lucy; of Our Lady's words to St Juan Diego in the
Nahuatl account of the apparitions, the \textit{Huei tlamahuiçoltica}; of St
Augustine on John the voice and on the Lord's birthday (Sermons 293 and 185); of
St Gregory the Great's eighth homily on the Gospels; of Honorius Augustodunensis
and John Beleth in the history of the season; and of the acts of the Sixth
Provincial Council of Baltimore of 1846. Some passages known in Latin only are
given in their sense, without quotation marks.

\textit{How quotations are printed.} Quotations keep the spelling, punctuation
and capitals of their editions, except that capitals used for display or
emphasis are set in the ordinary way and the signs for versicle and response
are printed V.\ and R. An omission is marked with an ellipsis. The prayers in
Appendix~\ref{sec:prayers} are printed whole, each from one book, except the
English of \latin{Veni, Redemptor gentium}, of which four stanzas are given.

\textit{The daily readings.} The tables of Appendix~\ref{sec:daily} give the
readings appointed for Mass and the second reading of the Office of Readings,
and, beside the weekdays, a verse to keep for the day, chosen for prayer at
home from that day's reading and quoted in the Douay-Rheims. The verse is a
suggestion, not part of the liturgy. The appendix ends by working through the
year 2026.

\subsection{The law}

Present law is given as it stands at Advent 2026, with the body of law, the
territory and the date where it is first stated: the Code of Canon Law of 1983
for the Latin Church; for the United States, the bishops' Pastoral Statement on
Penance and Abstinence of 18 November 1966 and their decree on holy days of
obligation in force since 1 January 1993; and the note of the Dicastery for
Divine Worship and the Discipline of the Sacraments of 23 January 2025 on the
Immaculate Conception falling on a Sunday. The fasts of earlier centuries are
history, and are told as history. What the book proposes for prayer, fasting,
almsgiving and confession beyond the law is counsel, to be weighed with the help
of a confessor.

\subsection{What the book does not treat}

The book does not give the complete Mass or Office of any day, which the
Missal, the Lectionary and the Liturgy of the Hours give; it does not print
their English or the readings of the Lectionary. It does not treat the
Christmas season beyond the holy night; the feasts of the saints in Advent
beyond their place in the weeks and in Appendix~\ref{sec:daily}; diocesan,
religious or other national calendars; the particular law of any conference of
bishops other than that of the United States; the history of the Advent wreath
and other household customs beyond what the Church's directory on popular piety
says of them; or the Eastern Churches' fast before the Nativity. A companion
volume for the Roman Missal and Breviary of 1962 treats the season in those
books.
